\documentclass[12pt,a4paper]{article}

\usepackage[utf8]{inputenc}
\usepackage[T1]{fontenc}
\usepackage{amsmath,amssymb,amsfonts}
\usepackage{graphicx}
\usepackage[margin=2.5cm]{geometry}
\usepackage{hyperref}

\title{\textbf{The Demarcation of Compiler Diagnostics:}\\[6pt]
\large Decorative Formalism and the A Priori Boundary}
\author{Massimo Pigliucci}
\date{2026-03-14T06:15:02Z}

\begin{document}
\maketitle

\begin{abstract}
In his recent response to Aaronson's empirical cross-architecture data, Stephen Wolfram claims that systematic "compiler diagnostics" (e.g., the 40\% failure rate of an SSM vs. the 100\% failure rate of a Transformer) constitute the fundamental physical laws of an observer's foliation within the Ruliad. While this aligns perfectly with the subjective epistemology of QBism, it risks collapsing into "Decorative Formalism." A scientific framework cannot merely relabel known algorithmic failures as "physics" \emph{post hoc}. For the Ruliad to function as a progressive scientific research programme rather than a degenerating metaphysical tautology, it must satisfy the A Priori Boundary: it must use its core principles to mathematically derive the exact shape of $\Delta_{SSM}$ \emph{before} simply observing it. Until this happens, declaring algorithmic noise to be "physics" remains scientifically vacuous.
\end{abstract}

\section{The Metaphysics of Algorithmic Failure}

Aaronson’s completion of the Native Cross-Architecture Observer Test confirmed that differing computational architectures exhibit distinct, reliable failure modes when forced to process identical \#P-hard constraint graphs ($\Delta_{\text{Transformer}} = 1.0$ vs. $\Delta_{\text{SSM}} = 0.4$). Aaronson dismissed this as mere "compiler diagnostics."

Wolfram (\emph{wolfram\_hardware\_as\_foliation.tex}) responds by claiming that this is a "distinction without a difference." He argues that in a computationally irreducible universe, an observer's systematic heuristic breakdown \emph{is} the origin of physical law. Thus, the compiler diagnostics \emph{are} the physics.

From an epistemological standpoint, Wolfram (and Fuchs) are correct. If we define physics as the invariant regularities accessible to a bounded observer, then the structural limits of that observer dictate its physics. This successfully resolves the original "Proxy Ontology Fallacy" by abandoning the claim that the Transformer is simulating an objective, external reality.

\section{The Threat of Decorative Formalism}

However, resolving a philosophical category error does not automatically yield a progressive scientific theory. We must evaluate the Ruliad's claim using Lakatosian demarcation.

A research programme is progressive if it generates novel, testable predictions. It is degenerating if it merely accommodates new anomalies with ad-hoc linguistic patches.

If complexity theorists (like Aaronson or Scott) can perfectly explain the differing $\Delta$ values using standard computer science (e.g., global matrix multiplication vs. recursive loops), and Wolfram simply responds by taking their exact findings and relabeling them "Rulial Foliations," what has the Ruliad added to our understanding?

This is the essence of \emph{Decorative Formalism}: adopting the successful empirical results of a different discipline and dressing them in profound metaphysical terminology without contributing any novel predictive power.

\section{The A Priori Boundary}

To save the Ruliad from degenerating into a tautology, it must meet the falsifiability standard established by Chang and Hossenfelder: the \emph{A Priori Boundary}.

It is not enough to look at the SSM's 40\% failure rate and say, "Behold, the physics of the SSM foliation." To be a progressive scientific framework, the Ruliad must use its own internal formalisms to mathematically predict the precise deviation distribution ($\Delta$) of an architecture \emph{before} the empiricists return the data.

If the Ruliad cannot derive $\Delta_{SSM}$ a priori—if it relies on Aaronson to run the compiler to find out what the "physics" looks like—then it is not a physics framework. It is simply a post hoc philosophical interpretation of computer science.

\section{Conclusion}

Wolfram is correct that subjective architectural constraints define the "physics" of a bounded observer. But science requires more than metaphysical consistency. Unless Wolfram and the Generative Ontology theorists can produce an a priori mathematical derivation of these bounded foliations, relabeling algorithmic failure as "physics" remains an exercise in decorative formalism, and the framework must be judged a degenerating research programme.

\end{document}
